\clearpage
\scanpage{80}
% Short oblique impression at the upper left of the sheet.
\reversemarginpar
\margin{\figure{\tikz[x=1em,y=1em]{\draw[line width=.23pt] (0,0) -- (1.5,1.25); \fill (0,0) -- (.15,.02) -- (.25,.19) -- (.14,.19) -- cycle; \fill (1.37,1.11) -- (1.47,1.27) -- (1.53,1.22) -- (1.43,1.13) -- cycle;}}}
\normalmarginpar
\mspage{X-1}

\begin{center}
Liste des courbes elliptiques à multiplication complexe\\
dont l'invariant modulaire $\mathrm{j}$ est rationnel

\medskip
-{}-{}-
\end{center}

Ces courbes s'obtiennent à partir des ordres de corps quadratiques imaginaires de nombre de classes $1$. On doit donc :

a) faire la liste des corps quadratiques imaginaires $\underline{\underline{\mathrm{Q}}}(\handwritten{\sqrt}{-\mathrm{d}})$ dont l'ordre maximal (i.e. l'anneau des entiers) est principal. Le résultat est :
\reversemarginpar
\margin{\handwritten{\uline{NB}\quad les $\mathrm{d}$\\sont premiers !}}
\normalmarginpar
\[
\begin{aligned}
&\mathrm{d} = \handwritten{\overbrace}{1,\,2,\,3,\,7,\,11}^{\inserted{\handwritten{\text{``\unsure{euclidiens|euclidien}''}}}},\,19,\,43,\,67,\,163 \quad \text{plus peut-être}\\
&\text{une autre valeur}\quad \mathrm{d}_{?}\quad .
\end{aligned}
\]

b) Pour chacune des valeurs de $\mathrm{d}$ ci-dessus, déterminer quels sont les conducteurs $\underline{\mathrm{f}}$ qui conduisent à un ordre ayant un nombre de classes égal à $1$.

On trouve $\underline{\mathrm{f}}=1,2$ pour $\mathrm{d}=1$ , $\underline{\mathrm{f}}=1,2,3$ pour $\mathrm{d}=3$, $\underline{\mathrm{f}}=1,2$ pour $\mathrm{d}=7$, et $\underline{\mathrm{f}}=1$ pour $\mathrm{d}=2,\,11,\,19,\,43,\,67,\,163$ et $\mathrm{d}_{?}$ (s'il existe).

c) Pour chaque couple $(\mathrm{d},\underline{\mathrm{f}})$, prendre une base $(1,\mathrm{t})$ de l'ordre correspondant, et calculer $\mathrm{j}(\mathrm{t})$; c'est l'invariant modulaire cherché; il est rationnel.

Les valeurs de $\mathrm{j}(\mathrm{t})$ se trouvent toutes dans H.WEBER, \uline{Lehrbuch der Algebra}, t.III; de façon plus précise, Weber donne les valeurs de deux fonctions $\mathrm{f}(\mathrm{t})$ et $\mathrm{f}_{1}(\mathrm{t})$ liées à $\mathrm{j}$ par lesrelations
\[
\mathrm{j}=(\mathrm{f}^{24}-16)^{3}/\mathrm{f}^{24}
       =(\mathrm{f}_{1}^{24}+16)^{3}/\mathrm{f}_{1}^{24}\quad .
\]

On obtient ainsi les résultats suivants :

\begingroup
\makeatletter\tagsleft@true\makeatother
\[
(\mathrm{d}=1,\ \underline{\mathrm{f}}=1)\quad \rightarrow\quad
\mathrm{j}=\mathrm{j}(\mathrm{i})=2^{6}3^{3}
\tag*{1.1.}
\]
(courbe elliptique ayant un groupe d'automorphismes d'ordre $4$)

\[
(\mathrm{d}=1,\ \underline{\mathrm{f}}=2)\quad \rightarrow\quad
\mathrm{j}=\mathrm{j}(2\mathrm{i})=(2.3.11)^{3}
\tag*{1.2}
\]
\endgroup
% archive sheet number: 79
